\scanpage{53}

% 1, cancelled
\deleted{Soit $\mathbb P\supset\mathbb K$ une \unsure{sur-catégorie|sous-catégorie} $k$-additive.}

% 2, cancelled upper pair
\diagram{\begin{tikzcd}[ampersand replacement=\&,execute at end picture={\draw (current bounding box.south west) -- (current bounding box.north east);}]
 \deleted{\underline{\operatorname{Hom}}_{k\text{-}\mathrm{lin}}(\mathbb P,B)}\arrow[r,shift left,"{\deleted{\rho}}"]\&\deleted{B^{\bullet}}\arrow[l,shift left,"{\deleted{\sigma}}"]
\end{tikzcd}}

% 2–4, upper right
\margin{\deleted{$B$ cat. additive avec $\Sigma$ \deleted{\unsure{de}} classe de suites exactes et $\Sigma$-noyaux}}

% 3, left, alongside triangle
\deleted{\unsure{pl. f.}}

% 3–4, cancelled triangle
\diagram{\begin{tikzcd}[ampersand replacement=\&,execute at end picture={\draw (current bounding box.south west) -- (current bounding box.north east);}]
 \deleted{\underline{\operatorname{Hom}}_{k\text{-}\mathrm{lin}}(\mathbb P,\widehat B)}\arrow[d]\&\\
 \deleted{\underline{\operatorname{Hom}}_{k\text{-}\mathrm{lin}}(\mathbb P,B)}\arrow[r,"{\deleted{\rho}}"]\&\deleted{B^{\bullet}}\arrow[ul,"{\deleted{\sigma}}"]
\end{tikzcd}}

% 5–6
$B$ catégorie additive, \deleted{avec noyaux de flèches,} \deleted{\inserted{classe $\Sigma$ de suites exactes}} où tt morphisme admet un $\Sigma$-noyau

% 7, new pair of functors
\diagram{\begin{tikzcd}[ampersand replacement=\&]
 \underline{\operatorname{Hom}}_{k,\unsure{\Sigma\text{-}\mathrm{sex}}}(\mathbb P_{\bullet},B)\arrow[r,shift left,"\rho"]\& B^{\bullet}\arrow[l,shift left,"{\unsure{\sigma|\substack{\sigma\\\mathrm{adj}}}}"]
\end{tikzcd}}

% 7, large note at right
% slanted
\margin{non, il faut condition de platitude ?}

% 8–11, below left Hom
cat. des foncteurs $k$-lin $F:\mathbb P\longrightarrow B$ qui transforment suites $\in\Sigma$ en suites exactes à \unsure{d|g}

% 12–13, defining brace
\[
 \left\{\begin{aligned}
 \rho(F)&=F(\mathbb E^{\bullet})\\
 \sigma(C^{\bullet})(L_{\bullet})&\simeq\operatorname{Hom}^{\bullet}_{k}(L_{\bullet}^{\vee},C^{\bullet})\left[=\operatorname{Hom}^{\bullet}_{k}(\operatorname{Hom}(L_{\bullet},\mathbb E^{\bullet}_{\deleted{\unsure{\bullet}}}),C^{\bullet})\right]
 \end{aligned}\right.
\]

% 13, below the second member
\margin{défini grâce à l'existence des $\Sigma$-noyaux}

% 14
Formule d'adjonction

% 15–16, with vertical identification
\diagram{\begin{tikzcd}[ampersand replacement=\&]
 \operatorname{Hom}(\rho(F),C^{\bullet})\arrow[r,"\sim"]\arrow[d,phantom,sloped,"\simeq"]\&\operatorname{Hom}(F,\sigma(C^{\bullet}))\\
 \operatorname{Hom}(F(\mathbb E^{\bullet}),C^{\bullet})\&
\end{tikzcd}}

% 17
Le deuxième mb est formé des hom fonctoriels en $L_{\bullet}$

% 18
\[
 F(L_{\bullet})\longrightarrow\operatorname{Hom}^{\bullet}_{k}(\underset{\substack{\Vert\\L^{\bullet}}}{L_{\bullet}^{\vee}},C^{\bullet})
\]

% 19
i.e. des hom ``fonctoriels en $L_{\bullet}$''

% 20
\[
 L^{\deleted{\unsure{k}}\bullet}\otimes_k F(L_{\bullet})\longrightarrow C^{\bullet}
\]

% 21
Or on a $L^i\simeq\operatorname{Hom}(L_{\bullet},\mathbb E^i)$

% 22
\[
 F(L)\otimes L^i\deleted{\otimes F(L_{\bullet})}\simeq\operatorname{Hom}(L_{\bullet},\mathbb E^i)\otimes F(L_{\bullet})\xrightarrow{\mathrm{taut.}}F(\mathbb E^i)
\]

% 23–24
et on trouve qu'on a ainsi la \deleted{``fonctorialité} compatibilité à l'op. diff. en $L^{\bullet}$ et $C^{\bullet}$, et ``fonctorialité

% archive, bottom left
% archive sheet number: 52
